% ---------------------------------------------------------
% TABELA EMENTA 45: Gestão de Pessoas e Relações no Trabalho
% ---------------------------------------------------------
\begin{xltabular}{\linewidth}{|X|p{2.5cm}|p{2.5cm}|}
\hline
\multirow{3}{=}{\textcolor{ifscgreen}{\textbf{Unidade Curricular:}}\\[3pt]\textbf{\large Gestão de Pessoas e Relações no Trabalho}} & \multicolumn{2}{p{\dimexpr5cm+2\tabcolsep+\arrayrulewidth\relax}|}{\textcolor{ifscgreen}{\textbf{Semestre:}} \textbf{5º e 6º (Ano 3)}} \\ \cline{2-3}
 & \textcolor{ifscgreen}{\textbf{CH EaD*:}} & \textcolor{ifscgreen}{\textbf{CH Total*:}} \\ \cline{2-3}
 & \textbf{00 h} & \textbf{80 h} \\ \hline
\endhead
\multicolumn{3}{|p{\dimexpr\linewidth-2\tabcolsep-2\arrayrulewidth\relax}|}{\textcolor{ifscgreen}{\textbf{Objetivos:}}} \\ \hline
\multicolumn{3}{|p{\dimexpr\linewidth-2\tabcolsep-2\arrayrulewidth\relax}|}{
\begin{itemize}[leftmargin=*,noitemsep,topsep=2pt]
 \item Compreender os fundamentos e os processos da gestão de pessoas nas organizações.
 \item Identificar práticas de recrutamento, seleção, desenvolvimento, avaliação e retenção de pessoas.
 \item Analisar a influência da liderança, da cultura organizacional e do clima no desempenho das equipes.
 \item Reconhecer a importância da ética, da diversidade, da inclusão e da qualidade de vida na gestão de pessoas.
 \item Aplicar conceitos básicos da gestão de pessoas em situações organizacionais.
\end{itemize}
} \\ \hline
\multicolumn{3}{|p{\dimexpr\linewidth-2\tabcolsep-2\arrayrulewidth\relax}|}{\textcolor{ifscgreen}{\textbf{Conteúdos:}}} \\ \hline
\multicolumn{3}{|p{\dimexpr\linewidth-2\tabcolsep-2\arrayrulewidth\relax}|}{
\begin{itemize}[leftmargin=*,noitemsep,topsep=2pt]
 \item Fundamentos e evolução da gestão de pessoas.
 \item Planejamento e dimensionamento da força de trabalho.
 \item Recrutamento, seleção e integração de pessoas.
 \item Desenvolvimento de pessoas, treinamento e educação corporativa.
 \item Avaliação de desempenho e desenvolvimento profissional.
 \item Liderança e gestão de equipes.
 \item Cultura e clima organizacional.
 \item Motivação e qualidade de vida no trabalho.
 \item Diversidade, inclusão e ética na gestão de pessoas.
 \item Legislação trabalhista aplicada à gestão de pessoas.
 \item Tendências contemporâneas da gestão de pessoas.
\end{itemize}
} \\ \hline
\multicolumn{3}{|p{\dimexpr\linewidth-2\tabcolsep-2\arrayrulewidth\relax}|}{\textcolor{ifscgreen}{\textbf{Estratégias de Ensino e Aprendizagem:}}} \\ \hline
\multicolumn{3}{|p{\dimexpr\linewidth-2\tabcolsep-2\arrayrulewidth\relax}|}{- **Metodologia:** Aulas expositivas dialogadas e metodologias ativas de aprendizagem, com enfoque teórico-prático, incluindo aprendizagem baseada em problemas, estudos de caso, oficinas, simulações e dinâmicas de grupo. Atividades podem incluir análise de documentos e práticas organizacionais, entrevistas, visitas técnicas e interações com organizações e profissionais da área. Utiliza sala de aula, laboratório de informática, sala multidisciplinar e espaços de organizações externas. Articula-se com Introdução à Administração, Organização e Processos, Gestão de Operações e Qualidade e Empreendedorismo, além de Língua Portuguesa, Sociologia, Sociedade e Trabalho e História.
- **Avaliação:** Contínua, processual e formativa, contemplando produções escritas, apresentações orais, resolução de estudos de caso, atividades práticas, relatórios, registros reflexivos e participação nas atividades propostas, de forma individual ou em grupo.} \\ \hline
\multicolumn{3}{|p{\dimexpr\linewidth-2\tabcolsep-2\arrayrulewidth\relax}|}{\textcolor{ifscgreen}{\textbf{Bibliografia Básica:}}} \\ \hline
\multicolumn{3}{|p{\dimexpr\linewidth-2\tabcolsep-2\arrayrulewidth\relax}|}{1. FIDELIS, Gilson José. \textbf{Gestão de pessoas: rotinas trabalhistas e dinâmicas do departamento de pessoal}. 6. ed. rev. atual. São Paulo: Érica, 2020.
2. FONTOURA, Iara P.; SABATOVSKI, Emilio. \textbf{Consolidação das leis trabalhistas}. 6. ed. Curitiba: Juruá, 2010.} \\ \hline
\multicolumn{3}{|p{\dimexpr\linewidth-2\tabcolsep-2\arrayrulewidth\relax}|}{\textcolor{ifscgreen}{\textbf{Bibliografia Complementar:}}} \\ \hline
\multicolumn{3}{|p{\dimexpr\linewidth-2\tabcolsep-2\arrayrulewidth\relax}|}{1. GIL, Antonio Carlos. \textbf{Gestão de pessoas: enfoque nos papéis profissionais}. São Paulo: Atlas, 2010.
2. IPONTELO, Juliana F.; CRUZ, Lucineide A. M. \textbf{Gestão de pessoas: manual de rotinas trabalhistas}. 8. ed. rev. e ampl. Brasília, DF: Senac-DF, 2016.
3. VERGARA, Sylvia Constant. \textbf{Gestão de pessoas}. 10. ed. São Paulo: Atlas, 2011.} \\ \hline
\end{xltabular}

\vspace{0.4cm}


% ---------------------------------------------------------
% TABELA EMENTA 46: Gestão Financeira
% ---------------------------------------------------------
\begin{xltabular}{\linewidth}{|X|p{2.5cm}|p{2.5cm}|}
\hline
\multirow{3}{=}{\textcolor{ifscgreen}{\textbf{Unidade Curricular:}}\\[3pt]\textbf{\large Gestão Financeira}} & \multicolumn{2}{p{\dimexpr5cm+2\tabcolsep+\arrayrulewidth\relax}|}{\textcolor{ifscgreen}{\textbf{Semestre:}} \textbf{5º e 6º (Ano 3)}} \\ \cline{2-3}
 & \textcolor{ifscgreen}{\textbf{CH EaD*:}} & \textcolor{ifscgreen}{\textbf{CH Total*:}} \\ \cline{2-3}
 & \textbf{00 h} & \textbf{80 h} \\ \hline
\endhead
\multicolumn{3}{|p{\dimexpr\linewidth-2\tabcolsep-2\arrayrulewidth\relax}|}{\textcolor{ifscgreen}{\textbf{Objetivos:}}} \\ \hline
\multicolumn{3}{|p{\dimexpr\linewidth-2\tabcolsep-2\arrayrulewidth\relax}|}{
\begin{itemize}[leftmargin=*,noitemsep,topsep=2pt]
 \item \textit{[Objetivos de aprendizagem pendentes de preenchimento pelo docente responsável]}
\end{itemize}
} \\ \hline
\multicolumn{3}{|p{\dimexpr\linewidth-2\tabcolsep-2\arrayrulewidth\relax}|}{\textcolor{ifscgreen}{\textbf{Conteúdos:}}} \\ \hline
\multicolumn{3}{|p{\dimexpr\linewidth-2\tabcolsep-2\arrayrulewidth\relax}|}{
\begin{itemize}[leftmargin=*,noitemsep,topsep=2pt]
 \item \textit{[Conteúdo programático pendente de preenchimento pelo docente responsável]}
\end{itemize}
} \\ \hline
\multicolumn{3}{|p{\dimexpr\linewidth-2\tabcolsep-2\arrayrulewidth\relax}|}{\textcolor{ifscgreen}{\textbf{Estratégias de Ensino e Aprendizagem:}}} \\ \hline
\multicolumn{3}{|p{\dimexpr\linewidth-2\tabcolsep-2\arrayrulewidth\relax}|}{\textit{[Metodologia de ensino e avaliação pendente de preenchimento pelo docente responsável]}} \\ \hline
\multicolumn{3}{|p{\dimexpr\linewidth-2\tabcolsep-2\arrayrulewidth\relax}|}{\textcolor{ifscgreen}{\textbf{Bibliografia Básica:}}} \\ \hline
\multicolumn{3}{|p{\dimexpr\linewidth-2\tabcolsep-2\arrayrulewidth\relax}|}{1. \textit{[1º título da Bibliografia Básica pendente de indicação docente e validação no Parecer da Biblioteca do Câmpus]}
2. \textit{[2º título da Bibliografia Básica pendente de indicação docente e validação no Parecer da Biblioteca do Câmpus]}} \\ \hline
\multicolumn{3}{|p{\dimexpr\linewidth-2\tabcolsep-2\arrayrulewidth\relax}|}{\textcolor{ifscgreen}{\textbf{Bibliografia Complementar:}}} \\ \hline
\multicolumn{3}{|p{\dimexpr\linewidth-2\tabcolsep-2\arrayrulewidth\relax}|}{1. \textit{[1º título da Bibliografia Complementar pendente de indicação docente e validação no Parecer da Biblioteca do Câmpus]}
2. \textit{[2º título da Bibliografia Complementar pendente de indicação docente e validação no Parecer da Biblioteca do Câmpus]}
3. \textit{[3º título da Bibliografia Complementar pendente de indicação docente e validação no Parecer da Biblioteca do Câmpus]}} \\ \hline
\end{xltabular}
